\documentclass{amsart}
% For dealing with references we use the comment environment
\usepackage{verbatim}
\newenvironment{reference}{\comment}{\endcomment}
%\newenvironment{reference}{}{}
\newenvironment{slogan}{\comment}{\endcomment}
\newenvironment{history}{\comment}{\endcomment}

% For commutative diagrams we use Xy-pic
\usepackage[all]{xy}

% We use 2cell for 2-commutative diagrams.
\xyoption{2cell}
\UseAllTwocells

% We use multicol for the list of chapters between chapters
\usepackage{multicol}

% This is generally recommended for better output
\usepackage{lmodern}
\usepackage[T1]{fontenc}

% For cross-file-references

% Package for hypertext links:
\usepackage{hyperref}

% For any local file, say "hello.tex" you want to link to please

% Theorem environments.
%
\theoremstyle{plain}
\newtheorem{theorem}[subsection]{Theorem}
\newtheorem{proposition}[subsection]{Proposition}
\newtheorem{lemma}[subsection]{Lemma}

\theoremstyle{definition}
\newtheorem{definition}[subsection]{Definition}
\newtheorem{example}[subsection]{Example}
\newtheorem{exercise}[subsection]{Exercise}
\newtheorem{situation}[subsection]{Situation}

\theoremstyle{remark}
\newtheorem{remark}[subsection]{Remark}
\newtheorem{remarks}[subsection]{Remarks}

\numberwithin{equation}{subsection}

% Macros
%
\def\lim{\mathop{\mathrm{lim}}\nolimits}
\def\colim{\mathop{\mathrm{colim}}\nolimits}
\def\Spec{\mathop{\mathrm{Spec}}}
\def\Hom{\mathop{\mathrm{Hom}}\nolimits}
\def\Ext{\mathop{\mathrm{Ext}}\nolimits}
\def\SheafHom{\mathop{\mathcal{H}\!\mathit{om}}\nolimits}
\def\SheafExt{\mathop{\mathcal{E}\!\mathit{xt}}\nolimits}
\def\Sch{\mathit{Sch}}
\def\Mor{\mathop{\mathrm{Mor}}\nolimits}
\def\Ob{\mathop{\mathrm{Ob}}\nolimits}
\def\Sh{\mathop{\mathit{Sh}}\nolimits}
\def\NL{\mathop{N\!L}\nolimits}
\def\CH{\mathop{\mathrm{CH}}\nolimits}
\def\proetale{{pro\text{-}\acute{e}tale}}
\def\etale{{\acute{e}tale}}
\def\QCoh{\mathit{QCoh}}
\def\Ker{\mathop{\mathrm{Ker}}}
\def\Im{\mathop{\mathrm{Im}}}
\def\Coker{\mathop{\mathrm{Coker}}}
\def\Coim{\mathop{\mathrm{Coim}}}

% Boxtimes
%
\DeclareMathSymbol{\boxtimes}{\mathbin}{AMSa}{"02}

%
% Macros for moduli stacks/spaces
%
\def\QCohstack{\mathcal{QC}\!\mathit{oh}}
\def\Cohstack{\mathcal{C}\!\mathit{oh}}
\def\Spacesstack{\mathcal{S}\!\mathit{paces}}
\def\Quotfunctor{\mathrm{Quot}}
\def\Hilbfunctor{\mathrm{Hilb}}
\def\Curvesstack{\mathcal{C}\!\mathit{urves}}
\def\Polarizedstack{\mathcal{P}\!\mathit{olarized}}
\def\Complexesstack{\mathcal{C}\!\mathit{omplexes}}
% \Pic is the operator that assigns to X its picard group, usage \Pic(X)
% \Picardstack_{X/B} denotes the Picard stack of X over B
% \Picardfunctor_{X/B} denotes the Picard functor of X over B
\def\Pic{\mathop{\mathrm{Pic}}\nolimits}
\def\Picardstack{\mathcal{P}\!\mathit{ic}}
\def\Picardfunctor{\mathrm{Pic}}
\def\Deformationcategory{\mathcal{D}\!\mathit{ef}}

\setlength{\emergencystretch}{3em}
\begin{document}
\title[Stacks Project, Tag 0ATX]{259 --- Canonical transitivity of twisted inverse image for separated finite-type morphisms}
\maketitle
\noindent Copyright (C) 2005--2025 Johan de Jong. Stacks Project authors. Source: \href{https://stacks.math.columbia.edu/tag/0ATX}{Tag 0ATX}.

\noindent Editorial difficulty note (not author text). Difficulty H2: The proof chooses compatible compactifications to construct the transitivity map, then dominates two such choices simultaneously. A five-square base-change argument establishes that the comparison agrees under these refinements. Thus the theorem proves a canonical composition law, not merely a formal adjoint identity..

\noindent Editorial provenance note (not author text). History: excerpt from revision \texttt{a0444\allowbreak 6e57e\allowbreak c1fbc\allowbreak 252a8\allowbreak 71afc\allowbreak ec775\allowbreak 2fb28\allowbreak 07b14}, duality.tex, lines 3943--4150. Reference presentation was adapted; original-author context and core support blocks were retained or added. The mathematical wording is unchanged. Licensed under GNU FDL 1.2 or later; no invariant sections or cover texts. See ../licenses/COPYING. Context blocks reproduce author definitions and supporting results; they are not additional counted problems.

\begin{situation}
\label{situation-shriek}
Here $S$ is a Noetherian scheme and $\textit{FTS}_S$ is the category whose
\begin{enumerate}
\item objects are schemes $X$ over $S$ such that the structure
morphism $X \to S$ is both separated and of finite type, and
\item morphisms $f : X \to Y$ between objects are morphisms
of schemes over $S$.
\end{enumerate}
\end{situation}

\begin{lemma}
\label{lemma-upper-shriek-composition}
In Situation \ref{situation-shriek} let $f : X \to Y$ and $g : Y \to Z$
be composable morphisms of $\textit{FTS}_S$. Then there is a canonical
isomorphism $(g \circ f)^! \to f^! \circ g^!$.
\end{lemma}

\begin{proof}
Choose a compactification $i : Y \to \overline{Y}$ of $Y$ over $Z$.
Choose a compactification $X \to \overline{X}$ of $X$ over
$\overline{Y}$. This uses More on Flatness, Theorem \href{https://stacks.math.columbia.edu/tag/0F41}{theorem nagata (Tag 0F41)}
and Lemma \href{https://stacks.math.columbia.edu/tag/0A9Z}{lemma compactifyable (Tag 0A9Z)} twice.
Let $\overline{a}$ be the
right adjoint of Lemma \href{https://stacks.math.columbia.edu/tag/0A9E}{lemma twisted inverse image (Tag 0A9E)} for
$\overline{X} \to \overline{Y}$ and let $\overline{b}$
be the
right adjoint of Lemma \href{https://stacks.math.columbia.edu/tag/0A9E}{lemma twisted inverse image (Tag 0A9E)} for
$\overline{Y} \to Z$.
Then $\overline{a} \circ \overline{b}$ is the
right adjoint of Lemma \href{https://stacks.math.columbia.edu/tag/0A9E}{lemma twisted inverse image (Tag 0A9E)} for
the composition $\overline{X} \to Z$.
Hence $g^! = i^* \circ \overline{b}$ and
$(g \circ f)^! = (X \to \overline{X})^* \circ \overline{a} \circ \overline{b}$.
Let $U$ be the inverse image of $Y$ in $\overline{X}$
so that we get the commutative diagram
$$
\xymatrix{
X \ar[r]_j \ar[d] & U \ar[dl] \ar[r]_{j'} & \overline{X} \ar[dl] \\
Y \ar[r]_i \ar[d] & \overline{Y} \ar[dl] \\
Z
}
$$
Let $\overline{a}'$ be the
right adjoint of Lemma \href{https://stacks.math.columbia.edu/tag/0A9E}{lemma twisted inverse image (Tag 0A9E)} for
$U \to Y$.
Then $f^! = j^* \circ \overline{a}'$. We obtain
$$
\gamma : (j')^* \circ \overline{a} \to \overline{a}' \circ i^*
$$
by (\href{https://stacks.math.columbia.edu/tag/0A9L}{equation sheafy (Tag 0A9L)}) and we can use it to define
$$
(g \circ f)^! =
(j' \circ j)^* \circ \overline{a} \circ \overline{b} =
j^* \circ (j')^* \circ \overline{a} \circ \overline{b}
\to
j^* \circ \overline{a}' \circ i^* \circ \overline{b} =
f^! \circ g^!
$$
which is an isomorphism on objects of $D_\QCoh^+(\mathcal{O}_Z)$ by
Lemma \href{https://stacks.math.columbia.edu/tag/0A9P}{lemma proper noetherian (Tag 0A9P)}. To finish the proof we show that
this isomorphism is independent of choices made.

\medskip\noindent
Suppose we have two diagrams
$$
\vcenter{
\xymatrix{
X \ar[r]_{j_1} \ar[d] & U_1 \ar[dl] \ar[r]_{j'_1} & \overline{X}_1 \ar[dl] \\
Y \ar[r]_{i_1} \ar[d] & \overline{Y}_1 \ar[dl] \\
Z
}
}
\quad\text{and}\quad
\vcenter{
\xymatrix{
X \ar[r]_{j_2} \ar[d] & U_2 \ar[dl] \ar[r]_{j'_2} & \overline{X}_2 \ar[dl] \\
Y \ar[r]_{i_2} \ar[d] & \overline{Y}_2 \ar[dl] \\
Z
}
}
$$
We can first choose a compactification $i : Y \to \overline{Y}$
with dense image of $Y$ over $Z$ which dominates both
$\overline{Y}_1$ and $\overline{Y}_2$,
see More on Flatness, Lemma \href{https://stacks.math.columbia.edu/tag/0ATU}{lemma compactifications cofiltered (Tag 0ATU)}.
By More on Flatness, Lemma \href{https://stacks.math.columbia.edu/tag/0ATV}{lemma right multiplicative system (Tag 0ATV)} and
Categories, Lemmas \href{https://stacks.math.columbia.edu/tag/04VI}{lemma morphisms right localization (Tag 04VI)} and
\href{https://stacks.math.columbia.edu/tag/04VJ}{lemma equality morphisms right localization (Tag 04VJ)}
we can choose a compactification $X \to \overline{X}$ with dense image of
$X$ over $\overline{Y}$ with morphisms $\overline{X} \to \overline{X}_1$
and $\overline{X} \to \overline{X}_2$ and such that the composition
$\overline{X} \to \overline{Y} \to \overline{Y}_1$ is equal to
the composition $\overline{X} \to \overline{X}_1 \to \overline{Y}_1$
and such that the composition
$\overline{X} \to \overline{Y} \to \overline{Y}_2$ is equal to
the composition $\overline{X} \to \overline{X}_2 \to \overline{Y}_2$.
Thus we see that it suffices to compare the maps
determined by our diagrams when we have a commutative diagram
as follows
$$
\xymatrix{
X \ar[rr]_{j_1} \ar@{=}[d] & &
U_1 \ar[d] \ar[ddll] \ar[rr]_{j'_1} & &
\overline{X}_1 \ar[d] \ar[ddll] \\
X \ar'[r][rr]^-{j_2} \ar[d] & &
U_2 \ar'[dl][ddll] \ar'[r][rr]^-{j'_2} & &
\overline{X}_2 \ar[ddll] \\
Y \ar[rr]^{i_1} \ar@{=}[d] & & \overline{Y}_1 \ar[d] \\
Y \ar[rr]^{i_2} \ar[d] & & \overline{Y}_2 \ar[dll] \\
Z
}
$$
and moreover the compactifications $X \to \overline{X}_1$ and
$Y \to \overline{Y}_2$ have dense image.
We use $\overline{a}_i$, $\overline{a}'_i$, $\overline{c}$, and
$\overline{c}'$ for the
right adjoint of Lemma \href{https://stacks.math.columbia.edu/tag/0A9E}{lemma twisted inverse image (Tag 0A9E)} for
$\overline{X}_i \to \overline{Y}_i$, $U_i \to Y$,
$\overline{X}_1 \to \overline{X}_2$, and $U_1 \to U_2$.
Each of the squares
$$
\xymatrix{
X \ar[r] \ar[d] \ar@{}[dr]|A & U_1 \ar[d] \\
X \ar[r] & U_2
}
\quad
\xymatrix{
U_2 \ar[r] \ar[d] \ar@{}[dr]|B & \overline{X}_2 \ar[d] \\
Y \ar[r] & \overline{Y}_2
}
\quad
\xymatrix{
U_1 \ar[r] \ar[d] \ar@{}[dr]|C & \overline{X}_1 \ar[d] \\
Y \ar[r] & \overline{Y}_1
}
\quad
\xymatrix{
Y \ar[r] \ar[d] \ar@{}[dr]|D & \overline{Y}_1 \ar[d] \\
Y \ar[r] & \overline{Y}_2
}
\quad
\xymatrix{
X \ar[r] \ar[d] \ar@{}[dr]|E & \overline{X}_1 \ar[d] \\
X \ar[r] & \overline{X}_2
}
$$
is cartesian (see
More on Flatness, Lemma \href{https://stacks.math.columbia.edu/tag/0ATU}{lemma compactifications cofiltered (Tag 0ATU)} part (c)
for A, D, E and recall that $U_i$ is the inverse image of $Y$
by $\overline{X}_i \to \overline{Y}_i$ for B, C) and hence
gives rise to a base change map (\href{https://stacks.math.columbia.edu/tag/0A9L}{equation sheafy (Tag 0A9L)}) as follows
$$
\begin{matrix}
\gamma_A : j_1^* \circ \overline{c}' \to j_2^* &
\gamma_B : (j_2')^* \circ \overline{a}_2 \to \overline{a}'_2 \circ i_2^* &
\gamma_C : (j_1')^* \circ \overline{a}_1 \to \overline{a}'_1 \circ i_1^* \\
\gamma_D : i_1^* \circ \overline{d} \to i_2^* &
\gamma_E : (j'_1 \circ j_1)^* \circ \overline{c} \to (j'_2 \circ j_2)^*
\end{matrix}
$$
Denote $f_1^! = j_1^* \circ \overline{a}'_1$,
$f_2^! = j_2^* \circ \overline{a}'_2$,
$g_1^! = i_1^* \circ \overline{b}_1$,
$g_2^! = i_2^* \circ \overline{b}_2$,
$(g \circ f)_1^! =
(j_1' \circ j_1)^* \circ \overline{a}_1 \circ \overline{b}_1$, and
$(g \circ f)^!_2 =
(j_2' \circ j_2)^* \circ \overline{a}_2 \circ \overline{b}_2$.
The construction given in the first paragraph of the proof
and in Lemma \href{https://stacks.math.columbia.edu/tag/0AA0}{lemma shriek well defined (Tag 0AA0)} uses
\begin{enumerate}
\item $\gamma_C$ for the map $(g \circ f)^!_1 \to f_1^! \circ g_1^!$,
\item $\gamma_B$ for the map $(g \circ f)^!_2 \to f_2^! \circ g_2^!$,
\item $\gamma_A$ for the map $f_1^! \to f_2^!$,
\item $\gamma_D$ for the map $g_1^! \to g_2^!$, and
\item $\gamma_E$ for the map $(g \circ f)^!_1 \to (g \circ f)^!_2$.
\end{enumerate}
We have to show that the diagram
$$
\xymatrix{
(g \circ f)^!_1 \ar[r]_{\gamma_E} \ar[d]_{\gamma_C} &
(g \circ f)^!_2 \ar[d]_{\gamma_B} \\
f_1^! \circ g_1^! \ar[r]^{\gamma_A \circ \gamma_D} & f_2^! \circ g_2^!
}
$$
is commutative. We will use
Lemmas \href{https://stacks.math.columbia.edu/tag/0ATQ}{lemma compose base change maps (Tag 0ATQ)} and
\href{https://stacks.math.columbia.edu/tag/0ATR}{lemma compose base change maps horizontal (Tag 0ATR)}
and with (abuse of) notation as in
Remark \href{https://stacks.math.columbia.edu/tag/0ATS}{remark going around (Tag 0ATS)} (in particular
dropping $\star$ products with identity transformations
from the notation).
We can write $\gamma_E = \gamma_A \circ \gamma_F$ where
$$
\xymatrix{
U_1 \ar[r] \ar[d] \ar@{}[rd]|F & \overline{X}_1 \ar[d] \\
U_2 \ar[r] & \overline{X}_2
}
$$
Thus we see that
$$
\gamma_B \circ \gamma_E = \gamma_B \circ \gamma_A  \circ \gamma_F
= \gamma_A \circ \gamma_B \circ \gamma_F
$$
the last equality because the two squares $A$ and $B$ only
intersect in one point (similar to the last argument in
Remark \href{https://stacks.math.columbia.edu/tag/0ATS}{remark going around (Tag 0ATS)}). Thus it suffices to prove that
$\gamma_D \circ \gamma_C = \gamma_B \circ \gamma_F$.
Since both of these are equal to the map (\href{https://stacks.math.columbia.edu/tag/0A9L}{equation sheafy (Tag 0A9L)})
for the square
$$
\xymatrix{
U_1 \ar[r] \ar[d] & \overline{X}_1 \ar[d] \\
Y \ar[r] & \overline{Y}_2
}
$$
we conclude.
\end{proof}
\end{document}
